\documentclass[english, parskip=full]{scrartcl}
\usepackage[utf8]{inputenc}  % Kodierung der Datei
\usepackage[ngerman]{babel}  % Sprache auf Deutsch 
\usepackage{color}
\usepackage{graphicx, gensymb}
\usepackage{amsfonts, cancel, amssymb}
\usepackage{amsmath, amsthm, array, esint}
\usepackage{booktabs, multirow}  % schönere Tabellen
\usepackage{caption}  % Anpassung der Captions
\usepackage{scrlayer-scrpage}    % Manipulation des Seitenstils
%\pagestyle{scrheadings}  % Stil für die Seite setzen
%\automark{section}  % Abschnittsnamen als Seitenbeschriftung 
%\setheadsepline{.5pt}    % Kopzeile durch Linie abtrennen
\usepackage[hidelinks]{hyperref}  % Links und weitere PDF-Features
\usepackage{typearea, tikz, pgffor, verbatim}
\usetikzlibrary{calc, quotes}
\usepackage[cache=false, outputdir=build]{minted}
\usepackage{placeins} % provides \FloatBarrier

\definecolor{bg}{rgb}{0.9,0.9,0.9}
\newcommand\numberthis{\addtocounter{equation}{1}\tag{\theequation}}
\newcommand{\bra}{}
\newcommand{\kom}{}
\newcommand{\brak}{}
\newcommand{\braket}{}
\newcommand{\intx}{}
\newcommand{\intr}{}
\newcommand{\kla}{}
\newcommand{\klam}{}
\newcommand{\klammer}{}
\newcommand{\oper}{}
\newcommand{\tecxt}{}
\newcommand{\Bra}{}
\newcommand{\Ket}{}
\renewcommand{\i}{\text{i}}
\newcommand{\e}{\text{e}}
\newcommand*{\Fourier}{\textsc{Fourier}\ }
\newcommand*{\F}{\mathcal{F}}
\newcommand*{\sinc}{\text{sinc\,}}
\newcommand{\conv}{\mathop{\scalebox{3.5}{\raisebox{-0.38ex}{\makebox[0.5\width][c]{$\ast$}}}}\limits}

\newcommand*{\titel}{02 Das Verfahren konjugierter Gradienten}
\newcommand*{\autor}{Joachim Schwardt und Julian Fleck}

\hypersetup{pdfauthor={\autor}, pdftitle={\titel}}

%\titlehead{titlehead}
\subject{Numerik II - Aufgaben}
\title{\titel}
%\author{\autor}
\author{\begin{tabular}{l|l}
Joachim Schwardt & 4768711 \\ 
Julian Fleck & 4759587\\ 
\end{tabular}}
\date{\today}

\begin{document}

%\begin{titlepage}
\maketitle  % Titel setzen
%\tableofcontents  % Inhaltsverzeichnis setzen
%\end{titlepage}

\section*{Aufgabe 1}
Gegeben sind
\begin{align}
A &:= \begin{pmatrix}
1 & 2 & 1 \\ 2 & 5 & 2 \\ 1 & 2 & 2
\end{pmatrix}&&\tecxt\text{und} & b &= \begin{pmatrix}
1 \\ 1 \\ 1 
\end{pmatrix}.
\end{align}
Offensichtlich ist $A^\top = A$.
Um zu zeigen, dass $A$ positiv definit ist, betrachten wir die Hauptminoren
\begin{align}
A_1 &= \begin{pmatrix}
1
\end{pmatrix}&&\tecxt\text{mit} & \det A_1 &= 1 > 0,\\
A_2 &= \begin{pmatrix}
1 & 2 \\ 2 & 5
\end{pmatrix}&&\tecxt\text{mit} & \det A_2 &= 5-4=1 > 0,\\
A_3 &= \begin{pmatrix}
1 & 2 & 1 \\ 2 & 5 & 2 \\ 1 & 2 & 2
\end{pmatrix}&&\tecxt\text{mit} & \det A_3 &= 10 + 4 + 4 - 5 - 4 - 8 = 1 > 0.
\end{align}
Nach dem Hauptminorenkriterium ist $A$ also positiv definit, da alle $\det A_{1,2,3} > 0$ sind.
\\
Gegeben ist der Startvektor $x^0 = (1,0,0)^\top$.
Der erste Iterationsschritt des CG-Verfahrens lautet dann
\begin{align}
r_0 &= b - Ax_0 = \begin{pmatrix}
1 \\ 1 \\ 1
\end{pmatrix} - \begin{pmatrix}
1 \\ 2 \\ 1
\end{pmatrix} = \begin{pmatrix}
0 \\ -1 \\ 0
\end{pmatrix},\\
d_0 &= r_0,\\
Ad_0 &= \begin{pmatrix}
-2 \\ -5 \\ -2
\end{pmatrix},\\
\alpha_0 &= \frac{r_0^\top r_0}{d_0^\top Ad_0} = \frac{1}{5},\\
x_1 &= x_0 + \alpha_0d_0 = \begin{pmatrix}
1 \\ -\frac{1}{5} \\ 0
\end{pmatrix}.
\end{align}
Der zweite Schritt lautet
\begin{align}
r_1 &= b - Ax_1 = r_0 - \alpha_0Ad_0 = \begin{pmatrix}
0 \\ -1 \\ 0
\end{pmatrix} - \frac{1}{5}\begin{pmatrix}
-2 \\ -5 \\ -2
\end{pmatrix} = \frac{1}{5}\begin{pmatrix}
2 \\ 0 \\ 2
\end{pmatrix},\\
\beta_1 &= \frac{r_1^\top r_1}{r_0^\top r_0} = \frac{\frac{4}{25} + \frac{4}{25}}{1} = \frac{8}{25},\\
d_1 &= r_1 + \beta_1d_0 = \frac{1}{5}\begin{pmatrix}
2 \\ 0 \\ 2
\end{pmatrix} + \frac{8}{25}\begin{pmatrix}
0 \\ -1 \\ 0
\end{pmatrix} = \frac{1}{25}\begin{pmatrix}
10 \\ -8 \\ 10
\end{pmatrix},\\
Ad_1 &= \frac{1}{25}\begin{pmatrix}
1 & 2 & 1 \\ 2 & 5 & 2 \\ 1 & 2 & 2
\end{pmatrix} \begin{pmatrix}
10 \\ -8 \\ 10
\end{pmatrix} = \frac{1}{25} \begin{pmatrix}
4 \\ 0 \\ 14
\end{pmatrix},\\
\alpha_1 &= \frac{r_1^\top r_1}{d_1^\top Ad_1} = \frac{\frac{8}{25}}{\frac{1}{25^2}(40 + 140)} = \frac{200}{180} = \frac{10}{9},\\
x_2 &= x_1 + \alpha_1d_1 = \begin{pmatrix}
1 \\ -\frac{1}{5} \\ 0
\end{pmatrix} + \frac{10}{9}\frac{1}{25} \begin{pmatrix}
10 \\ -8 \\ 10
\end{pmatrix} = \\
&= \frac{1}{45}\begin{pmatrix}
45 \\ -9 \\ 0
\end{pmatrix} + \frac{2}{45}\begin{pmatrix}
10 \\ -8 \\ 10
\end{pmatrix} = \frac{1}{45}\begin{pmatrix}
65 \\ -25 \\ 20
\end{pmatrix} = \frac{1}{9}\begin{pmatrix}
13 \\ -5 \\ 4
\end{pmatrix}.
\end{align}
Der dritte Schritt lautet
\begin{align}
r_2 &= b - Ax_2 = r_1 - \alpha_1Ad_1 = \frac{1}{5}\begin{pmatrix}
2 \\ 0 \\ 2
\end{pmatrix} - \frac{10}{9}\frac{1}{25} \begin{pmatrix}
4 \\ 0 \\ 14
\end{pmatrix} = \\
&= \frac{1}{45}\begin{pmatrix}
10 \\ 0 \\ -10
\end{pmatrix} = \frac{1}{9}\begin{pmatrix}
2 \\ 0 \\ -2
\end{pmatrix},\\
\beta_2 &= \frac{r_2^\top r_2}{r_1^\top r_1} = \frac{\frac{4}{81} + \frac{4}{81}}{\frac{8}{25}} = \frac{25}{81},\\
d_2 &= r_2 + \beta_2d_1 = \frac{1}{9}\begin{pmatrix}
2 \\ 0 \\ -2
\end{pmatrix} + \frac{25}{81} \frac{1}{25}\begin{pmatrix}
10 \\ -8 \\ 10
\end{pmatrix} = \frac{1}{81}\begin{pmatrix}
28 \\ -8 \\ -8
\end{pmatrix},\\
Ad_2 &= \frac{1}{81}\begin{pmatrix}
1 & 2 & 1 \\ 2 & 5 & 2 \\ 1 & 2 & 2
\end{pmatrix} \begin{pmatrix}
28 \\ -8 \\ -8
\end{pmatrix} = \frac{1}{81} \begin{pmatrix}
4 \\ 0 \\ -4
\end{pmatrix},\\
\alpha_2 &= \frac{r_2^\top r_2}{d_2^\top Ad_2} = \frac{\frac{8}{81}}{\frac{4}{81^2}(28 + 8)} = \frac{81}{18} = \frac{9}{2},\\
x_3 &= x_2 + \alpha_2d_2 = \frac{1}{9}\begin{pmatrix}
13 \\ -5 \\ 4
\end{pmatrix} + \frac{9}{2} \frac{1}{81} \begin{pmatrix}
28 \\ -8 \\ -8
\end{pmatrix} = \frac{1}{18}\begin{pmatrix}
26 + 28 \\ -10 - 8 \\ 8 - 8
\end{pmatrix} = \begin{pmatrix}
3 \\ -1 \\ 0
\end{pmatrix}.
\end{align}
Da $Ax_3 = (3-2, 6-5, 3-2)^\top = b$, ist $x^3=x^\ast$ die Lösung des Gleichungssystems. 
Das Verfahren bricht ab, weil dementsprechend $r^3 = 0$.

%import numpy as np
%
%A = np.array([[1,2,1],[2,5,2],[1,2,2]])
%b = np.array([1,1,1])
%x0 = np.array([1,0,0])

%def cg(A, b, x0=0, N=4):
%    r0 = b - A @ x0
%    d = r0
%    r = r0
%    for i in range(N):
%        Ad = A @ d
%        alpha = (r @ r) / (d @ Ad)
%        x0 = x0 + alpha * d
%        r = r0 - alpha * Ad
%        beta = (r @ r) / (r0 @ r0)
%        d = r + beta * d
%        print(f"{Ad = } \n{alpha = } \n{x0 = } \n{r = } \n{beta = } \n{d = } \n{r0 = }\n\n")
%        r0 = r
%    return x0

%Berechnung des nullten Residuums $r_0 = b - Ax_{start}$
%Dies ergibt: $[\frac{0}{1} , \frac{-1}{1} , \frac{0}{1}]$
%Dieser Vektor wird auch als erste Suchrichtung verwendet.
%
%Mit der Suchrichtung wird die neue Schrittweite $\alpha_0 = \frac{r_0^Tr_0}{d^T_0Ad_0^T} = \frac{1}{5}$
%Damit ergibt sich die neue iterierte Lösung des Gleichungssystems $x_1 = x_0 + \alpha_0 d_0 = [\frac{1}{1} , \frac{0}{1} , \frac{0}{1}] + \frac{1}{5} \cdot [\frac{0}{1} , \frac{-1}{1} , \frac{0}{1}] = [\frac{1}{1} , \frac{-1}{5} , \frac{0}{1}]$
%Nun kann das neue Residuum $r_1 = b - A x_1 = [\frac{2}{5} , \frac{0}{1} , \frac{2}{5}]$ errechnet werden.
%Aus dem neuen Residuum wird nun das Gram-Schmidt-Verfahren durchgeführt um die neue Suchrichtung zu bestimmen.
%Die neue Suchrichtung $d_1 = r_1 + \beta_{(1)} d_0$ ergibt sich mit $\beta_{(1)} = \frac{8}{25}$ zu $d_1 = [\frac{2}{5} , \frac{-8}{25} , \frac{2}{5}]$
%
%
%Mit der Suchrichtung wird die neue Schrittweite $\alpha_1 = \frac{r_1^Tr_1}{d^T_1Ad_1^T} = \frac{10}{9}$
%Damit ergibt sich die neue iterierte Lösung des Gleichungssystems $x_2 = x_1 + \alpha_1 d_1 = [\frac{1}{1} , \frac{-1}{5} , \frac{0}{1}] + \frac{10}{9} \cdot [\frac{2}{5} , \frac{-8}{25} , \frac{2}{5}] = [\frac{13}{9} , \frac{-5}{9} , \frac{4}{9}]$
%Nun kann das neue Residuum $r_2 = b - A x_2 = [\frac{2}{9} , \frac{0}{1} , \frac{-2}{9}]$ errechnet werden.
%Aus dem neuen Residuum wird nun das Gram-Schmidt-Verfahren durchgeführt um die neue Suchrichtung zu bestimmen.
%Die neue Suchrichtung $d_2 = r_2 + \beta_{(2)} d_1$ ergibt sich mit $\beta_{(2)} = \frac{25}{81}$ zu $d_2 = [\frac{28}{81} , \frac{-8}{81} , \frac{-8}{81}]$
%
%
%Mit der Suchrichtung wird die neue Schrittweite $\alpha_2 = \frac{r_2^Tr_2}{d^T_2Ad_2^T} = \frac{9}{2}$
%Damit ergibt sich die neue iterierte Lösung des Gleichungssystems $x_3 = x_2 + \alpha_2 d_2 = [\frac{13}{9} , \frac{-5}{9} , \frac{4}{9}] + \frac{9}{2} \cdot [\frac{28}{81} , \frac{-8}{81} , \frac{-8}{81}] = [\frac{3}{1} , \frac{-1}{1} , \frac{0}{1}]$
%Nun kann das neue Residuum $r_3 = b - A x_3 = [\frac{0}{1} , \frac{0}{1} , \frac{0}{1}]$ errechnet werden.
%Aus dem neuen Residuum wird nun das Gram-Schmidt-Verfahren durchgeführt um die neue Suchrichtung zu bestimmen.
%Die neue Suchrichtung $d_3 = r_3 + \beta_{(3)} d_2$ ergibt sich mit $\beta_{(3)} = \frac{0}{1}$ zu $d_3 = [\frac{0}{1} , \frac{0}{1} , \frac{0}{1}]$
%
%
%Das neue Residuum und die neue Suchrichtung sind der Nullvektor, somit bricht das Verfahren ab und wir haben die Lösung des Gleichungssystems $x^* = [\frac{3}{1} , \frac{-1}{1} , \frac{0}{1}]$.


\newpage
\section*{Aufgabe 2}
Sei $A\in\mathbb{R}^{n\times n}$ symmetrisch und positiv definit.
Gegeben ist das lineare Gleichungssystem $Ax=b$, wobei $b=c_1v_1 + c_2v_2$ und $v_{1,2}$ Eigenvektoren von $A$ zum Eigenwert $\lambda_{1,2}$ sind.
Insbesondere sind die Eigenvektoren orthonormal, also $v_j^\top v_k = \delta_{jk}$.
Zu zeigen ist, dass das CG-Verfahren nach zwei Iterationsschritten abbricht.
Definiere dazu die Residuen $r_k := b- Ax_k$, wobei aufgrund des Startwertes $x_0 = 0$ einfach $r_0 = b$ ist.
Für $c_1=c_2=0$ wäre $b=0$ und die Lösung damit trivialerweise bereits gleich dem Startpunkt.
Sei im Folgenden also mindestens ein Koeffizient $c_{1,2} \neq 0$.
Der erste Iterationsschritt lautet dann
\begin{align}
d_0 &= r_0 = b,\\
\alpha_0 &= \frac{r_0^\top r_0}{d_0^\top Ad_0} = \frac{b^\top b}{b^\top Ab} = \frac{c_1^2 + c_2^2}{c_1^2\lambda_1 + c_2^2\lambda_2},\\
x_1 &= x_0 + \alpha_0 d_0 = \frac{c_1^2 + c_2^2}{c_1^2\lambda_1 + c_2^2\lambda_2}(c_1v_1 + c_2v_2),\\
r_1 &= b - Ax_1 = (c_1v_1 + c_2v_2) - \frac{c_1^2 + c_2^2}{c_1^2\lambda_1 + c_2^2\lambda_2}(c_1\lambda_1 v_1 + c_2\lambda_2 v_2) = \\
&= c_1v_1 \kla\left( 1 - \frac{c_1^2\lambda_1 + c_2^2\lambda_1}{c_1^2\lambda_1 + c_2^2\lambda_2} \right) + c_2v_2 \kla\left( 1 - \frac{c_1^2\lambda_2 + c_2^2\lambda_2}{c_1^2\lambda_1 + c_2^2\lambda_2} \right) =\\
&= c_1v_1 \frac{c_2^2(\lambda_2 - \lambda_1)}{c_1^2\lambda_1 + c_2^2\lambda_2} + c_2v_2 \frac{c_1^2(\lambda_1 - \lambda_2)}{c_1^2\lambda_1 + c_2^2\lambda_2},
\end{align}
wobei das Residuum noch zu 
\begin{align}
r_1 &= \frac{c_1c_2 (\lambda_2 - \lambda_1)}{c_1^2\lambda_1 + c_2^2\lambda_2}\kla\left( c_2v_1 - c_1v_2 \right)
\end{align}
vereinfacht werden kann.
Für $\lambda_1 = \lambda_2$ oder $c_1=0$ oder $c_2 = 0$ ist $r_1=0$, womit wir bereits fertig wären.
Sei also im Folgenden $\lambda_1\neq \lambda_2$ und $c_{1,2}\neq 0$.
\\
Explizit den zweiten Iterationsschritt durchzuführen hat sich als eine enorm aufwendige Aufgabe erwiesen, weshalb wir hier einen anderen Ansatz wählen.
Aus der Vorlesung ist bekannt, dass die Residuen die Relation
\begin{align}
r_j^\top r_k &= 0
\end{align}
für alle $k<j$ erfüllen (also orthogonal sind).
Das CG-Verfahren wird auch in allen weiteren Schritten immer nur Vektoren liefern, die Linearkombinationen der $v_{1,2}$ sind.
Daher können wir zunächst $r_2 =: \gamma_1c_1v_1 + \gamma_2c_2v_2$ für zwei Parameter $\gamma_{1,2}$ definieren.
Die Orthogonalitätsrelationen liefern dann
\begin{align}
r_2^\top r_1 &= \kla\left( \gamma_1c_1v_1 + \gamma_2c_2v_2 \right)^\top  \frac{c_1c_2 (\lambda_2 - \lambda_1)}{c_1^2\lambda_1 + c_2^2\lambda_2} \kla\left( c_2v_1 - c_1v_2 \right) = \\
&= \frac{c_1^2c_2^2 (\lambda_2 - \lambda_1)}{c_1^2\lambda_1 + c_2^2\lambda_2}\kla\left( \gamma_1 - \gamma_2 \right) \overset{!}{=} 0, \label{eqn:r2 r1 orthogonal}\\
r_2^\top r_0 &= \kla\left( \gamma_1c_1v_1 + \gamma_2c_2v_2 \right)^\top  \kla\left( c_1v_1+c_2v_2 \right) = \\
&= \gamma_1c_1^2 + \gamma_2c_2^2 \overset{!}{=} 0. \label{eqn:r2 r0 orthogonal}
\end{align}
Da $c_{1,2}\neq 0 $ folgt aus \autoref{eqn:r2 r1 orthogonal} zunächst $\gamma_1=\gamma_2$.
Eingesetzt in \autoref{eqn:r2 r0 orthogonal} folgt daraus $\gamma_1=0$ und somit $r_2 = 0$.

\newpage
\section*{Aufgabe 3}
Sei $A\in\mathbb{R}^{n\times n}$ und $r\in\mathbb{R}^n$. 
Dann ist der $k$-te \textsc{Krylov}-Raum zu $r$ und $A$ für $k\ge 0$ definiert durch
\begin{align}
\mathcal{D}_k(r, A) &:= \tecxt\text{span\,}\klammer\left\{ r, Ar, \dots, A^{k-1}r \right\}.
\end{align}
Sei $A$ zusätzlich regulär, $b\in\mathbb{R}^n$ und $x^\ast = A^{-1}b$.
Für $x_k\in\mathbb{R}^n$ definiere $r_k := b-Ax_k$, wobei insbesondere $r_0\neq 0$ gefordert sei.
Definiere noch $e_k := x^\ast - x_k$, woraus $Ae_k = b - Ax_k = r_k$ folgt.
\\
Angenommen es gilt $\mathcal{D}_{k+1}(r_0, A) = \mathcal{D}_k(r_0, A)$ für ein $k\in\mathbb{N}$.
Dann folgt aus $A^kr_0 \in \mathcal{D}_k(r_0, A)$ insbesondere, dass es $c_0,\dots,c_{k-1}\in\mathbb{R}$ gibt mit
\begin{align}
A^k r_0 &= \sum_{j=0}^{k-1} c_jA^jr_0 = \sum_{j=0}^{k-1} c_jA^{j+1}e_0. \label{eqn:Ak r0 Linearkombination}
\end{align}
Wir wollen diese Gleichung so umformen, dass wir eine Darstellung für $e_0$ erhalten.
Für $c_0\neq 0 $ gilt dabei einfach
\begin{align}
c_0Ae_0 &= A^kr_0 - \sum_{j=1}^{k-1}c_jA^{j+1}e_0,\\
\Rightarrow\ \ \ e_0 &= \frac{1}{c_0}\kla\left( A^{k-1}r_0 - \sum_{j=1}^{k-1}c_jA^{j-1}r_0 \right).
\end{align}
Da dieser Ausdruck eine Linearkombination der $r_0, Ar_0, \dots, A^{k-1}r_0$ ist, folgt $e_0 \in \mathcal{D}_k(r_0, A)$, was äquivalent zu $x^\ast \in x^0 + \mathcal{D}_k(r_0, A)$ ist.
\\
Für $c_0 = 0$ können wir \autoref{eqn:Ak r0 Linearkombination} auch nach $e_0$ umstellen.
Wir nehmen noch an, dass die ersten $q$ Koeffizienten gleich Null sind, also $c_0 = \dots = c_{q-1} = 0$ für $q\in\{1,\dots,k-1\}$.
Dann gilt
\begin{align}
A^kr_0 &= \sum_{j=q}^{k-1}c_jA^{j+1}e_0 = c_qA^{q+1}e_0 + \sum_{j=q+1}^{k-1}c_jA^jr_0,\\
\Rightarrow\ \ \ e_0 &= \frac{1}{c_q}\kla\left( A^{k-q-1}r_0 - \sum_{j=q+1}^{k-1}c_jA^{j-q-1}r_0 \right),
\end{align}
was ebenfalls eine Linearkombination in $\mathcal{D}_k(r_0, A)$ ist.
\\
Angenommen es gilt $e_0 \in \mathcal{D}_k(r_0, A)$.
Dann gibt es $c_0,\dots,c_{k-1}\in\mathbb{R}$, sodass
\begin{align}
e_0 &= \sum_{j=0}^{k-1}c_jA^jr_0. \label{eqn:e0 als Linearkombination von Aj r0}
\end{align}
Für $c_{k-1} \neq 0$ führt das Multiplizieren mit $A$ dann zu 
\begin{align}
r_0 &= Ae_0 = \sum_{j=0}^{k-1}c_jA^{j+1}r_0 = \sum_{j=1}^{k}c_{j-1} A^jr_0,\\
A^kr_0 &= \frac{1}{c_{k-1}}\kla\left( r_0 - \sum_{j=1}^{k-1}c_{j-1}A^jr_0 \right).
\end{align}
Das ist eine Linearkombination der $r_0, \dots, A^{k-1}r_0$, womit $A^kr_0 \in \mathcal{D}_k(r_0, A)$.
Insbesondere gilt dann $\mathcal{D}_{k+1}(r_0, A) = \mathcal{D}_k(r_0, A)$, da der neue Eintrag linear abhängig von den anderen Vektoren ist.
\\
Für $c_{k-1} = 0$ führen wir wieder $q\in\{1,\dots,k-1\}$ ein, wobei diesmal $c_{k-1} = \dots = c_{k-q} = 0$ sei.
Dann können wir \autoref{eqn:e0 als Linearkombination von Aj r0} durch Multiplizieren mit $A^{q+1}$ wieder nach $A^kr_0$ umstellen, da
\begin{align}
A^{q+1}e_0 &= \sum_{j=0}^{k-q-1}c_jA^{j+q+1}r_0 = c_{k-q-1}A^kr_0 + \sum_{j=0}^{k-q-2}c_jA^{j+q+1}r_0,\\
\Rightarrow\ \ \ A^kr_0 &= \frac{1}{c_{k-q-1}} \kla\left( A^{q-1}r_0 - \sum_{j=0}^{k-q-2}c_jA^{j+q+1}r_0 \right),
\end{align}
was ebenfalls eine Linearkombination in $\mathcal{D}_k(r_0, A)$ ist.


%\begin{thebibliography}{9}
%\bibitem{example} description, \url{https://www.exmapleurl}, day.\,month~2021.
%\end{thebibliography}

\end{document}
